\documentclass[11pt,a4paper]{article}
\usepackage[utf8]{inputenc}
\usepackage[spanish]{babel}
\usepackage{amsmath,amsfonts,amssymb}
\usepackage{geometry}
\usepackage{listings}
\usepackage{xcolor}
\usepackage{booktabs}

\geometry{margin=1in}

\definecolor{codegreen}{rgb}{0,0.6,0}
\definecolor{backcolour}{rgb}{0.95,0.95,0.92}

\lstset{
    backgroundcolor=\color{backcolour},
    commentstyle=\color{codegreen},
    keywordstyle=\color{blue},
    basicstyle=\ttfamily\footnotesize,
    breaklines=true,
    numbers=left,
    showstringspaces=false,
}

\title{Simulación de Confiabilidad: Selección de un Intervalo Preventivo en Mantenimiento Industrial}
\author{Cátedra de Gestión de Activos - MIT}
\date{2026}

\begin{document}

\maketitle

%=========================================================
\section{Contexto: el dilema del gestor de mantenimiento}
%=========================================================
Una línea de ensamblaje automatizada depende de un motor crítico. Cuando el motor falla de manera inesperada
(\textit{mantenimiento correctivo}), la línea se detiene y la empresa incurre en costos de parada y de reemplazo.
Alternativamente, el motor puede reemplazarse de forma programada (\textit{mantenimiento preventivo}), evitando fallas pero
asumiendo un costo fijo y realizando el cambio antes de agotar totalmente la vida útil del equipo.

En este ejercicio nos enfocaremos en una política sencilla y muy usada en confiabilidad:
\[
\textbf{Reemplazo por edad:}\quad \text{reemplazar preventivamente al cumplir } T_p \text{ horas si no ha fallado antes.}
\]
La pregunta es: \textbf{¿qué valor de $T_p$ minimiza el costo esperado por hora de operación?}

\subsection*{Costos y tiempos (datos del problema)}
\begin{itemize}
    \item \textbf{Preventivo:} costo fijo $C_p = \$1{,}200$ y se realiza en una hora planificada (sin costo de parada).
    \item \textbf{Correctivo:} si ocurre una falla, el reemplazo cuesta $C_r = \$1{,}200$ \emph{y} hay un costo de parada de
    \$5{,}000 por hora durante el tiempo de reparación. Denotemos por $D_c$ el \textbf{tiempo de reparación} (en horas) bajo correctivo.
\end{itemize}

\noindent\textit{Nota:} en tu versión original aparecía ``\$5,000 por hora'' pero luego usabas $C_c=\$5,000$ como si fuese un costo fijo.
Para que el modelo sea consistente, aquí separamos:
\[
\text{costo correctivo} = C_r + (\text{costo por hora de paro})\times D_c.
\]
Si no quieres modelar $D_c$, puedes fijarlo como constante (p.ej.\ $D_c=1$ hora) y ya.

%=========================================================
\section{Modelo aleatorio: tiempo a la falla}
%=========================================================
Sea $TTF$ el \textbf{tiempo a la falla} (\textit{time to failure}). Supondremos:
\[
TTF \sim \text{Weibull}(k,\lambda), \qquad k=2.5,\;\lambda=1200\;\text{horas}.
\]
El parámetro $k>1$ captura un \emph{riesgo creciente} (desgaste). La función de confiabilidad es:
\[
R(t)=\mathbb{P}(TTF>t)=\exp\!\left(-\left(\frac{t}{\lambda}\right)^k\right).
\]

%=========================================================
\section{Política de decisión: reemplazo por edad}
%=========================================================
Dado un intervalo preventivo $T_p$:
\begin{itemize}
    \item Si $TTF<T_p$, ocurre falla antes del preventivo $\Rightarrow$ correctivo.
    \item Si $TTF\ge T_p$, llega vivo al preventivo $\Rightarrow$ preventivo.
\end{itemize}

Cada ciclo de operación termina cuando el motor falla o cuando se hace el preventivo. Por tanto, el \textbf{tiempo de operación del ciclo} es:
\[
\tau(T_p) = \min\{TTF,\,T_p\}.
\]

%=========================================================
\section{¿Cuál es la métrica correcta? Costo promedio por unidad de tiempo}
%=========================================================
El objetivo natural (y estándar) en reemplazo por edad es minimizar el \textbf{costo de largo plazo por hora}, que se estima como:
\[
\widehat{g}(T_p) \approx \frac{\sum_{k=1}^{N} C^{(k)}(T_p)}{\sum_{k=1}^{N} \tau^{(k)}(T_p)}.
\]
Aquí $C^{(k)}(T_p)$ es el costo del ciclo $k$ y $\tau^{(k)}(T_p)$ su duración.
Esta razón es el estimador Monte Carlo del costo promedio por unidad de tiempo (razón de promedios).

\noindent\textit{Nota:} en tu texto original se definía $CPH=\text{Costo}/\text{Tiempo}$ por ciclo y luego se promediaba.
Eso produce $\mathbb{E}[C/\tau]$, que en general \emph{no} es igual a $\mathbb{E}[C]/\mathbb{E}[\tau]$.
Para costo promedio de largo plazo (tipo \textit{renewal-reward}), la razón de sumas es el objeto correcto.

%=========================================================
\section{Algoritmo de simulación (paso a paso)}
%=========================================================
Para un valor fijo de $T_p$:
\begin{enumerate}
    \item Simular un $TTF$ de la Weibull$(k,\lambda)$.
    \item Determinar el evento:
    \begin{itemize}
        \item Si $TTF<T_p$: el ciclo termina por falla (correctivo).
        \item Si $TTF\ge T_p$: el ciclo termina por reemplazo preventivo.
    \end{itemize}
    \item Calcular:
    \[
    \tau = \min\{TTF,T_p\},
    \]
    y el costo del ciclo:
    \[
    C =
    \begin{cases}
      C_p, & \text{si } TTF\ge T_p,\\[4pt]
      C_r + 5000\cdot D_c, & \text{si } TTF<T_p.
    \end{cases}
    \]
    \item Repetir $N$ ciclos y estimar $g(T_p)=\dfrac{\sum C}{\sum \tau}$.
\end{enumerate}

%=========================================================
\section{Métricas adicionales (opcional, para discusión)}
%=========================================================
Además del costo promedio por hora, dos métricas útiles son:
\begin{itemize}
    \item \textbf{Confiabilidad al preventivo:} $R(T_p)=\mathbb{P}(TTF\ge T_p)$.
    \item \textbf{Frecuencia de correctivos:} proporción de ciclos con $TTF<T_p$ (aprox. $\widehat{\mathbb{P}}(TTF<T_p)$).
\end{itemize}
Si se modela $D_c$ como tiempo de reparación (y también un tiempo preventivo $D_p$), entonces se puede estimar disponibilidad.
Si el preventivo es ``en hora planificada'' y no descuenta disponibilidad operativa, basta con fijar $D_p=0$ en ese cálculo.

%=========================================================
\section{Implementación en Python (nivel principiante)}
%=========================================================
\begin{lstlisting}[language=Python, caption=Simulacion de reemplazo por edad y costo promedio por hora]
import numpy as np

# -----------------------
# 1) Parametros del modelo
# -----------------------
k_weibull = 2.5
lambda_weibull = 1200.0  # horas

C_p = 1200.0            # costo preventivo
C_r = 1200.0            # costo de reemplazo (tambien en correctivo)
costo_parada_hora = 5000.0

D_c = 1.0               # horas de reparacion en correctivo (simplificacion)
N = 10000               # numero de ciclos simulados

np.random.seed(2026)

# -----------------------
# 2) Funcion para simular un TTF Weibull(k, lambda)
#    Nota: numpy.random.weibull(k) devuelve una Weibull(k) con escala 1
#    Para escala lambda: TTF = lambda * W
# -----------------------
def simular_ttf():
    w = np.random.weibull(k_weibull)
    return lambda_weibull * w

# -----------------------
# 3) Evaluar una politica T_p
#    Estimador correcto: (sum costos) / (sum tiempos)
# -----------------------
def evaluar_politica(T_p):
    suma_costos = 0.0
    suma_tiempos = 0.0
    n_correctivos = 0

    for i in range(N):
        ttf = simular_ttf()

        # Duracion del ciclo
        tau = min(ttf, T_p)
        suma_tiempos += tau

        # Costo del ciclo
        if ttf < T_p:
            # Correctivo
            n_correctivos += 1
            costo = C_r + costo_parada_hora * D_c
        else:
            # Preventivo
            costo = C_p

        suma_costos += costo

    costo_promedio_por_hora = suma_costos / suma_tiempos
    fraccion_correctivos = n_correctivos / N
    return costo_promedio_por_hora, fraccion_correctivos

# -----------------------
# 4) Probar varios intervalos
# -----------------------
for T_p in [600, 800, 1000, 1200, 1400]:
    cph, frac_c = evaluar_politica(T_p)
    print(f"T_p = {T_p:4d} h | costo/hora = ${cph:6.2f} | frac. correctivos = {frac_c:5.2%}")
\end{lstlisting}

\end{document}
